\section*{Заключение}
\addcontentsline{toc}{section}{Заключение}
\label{sec:conclusion}

В ходе выполнения курсового проекта было разработано веб-приложение LogicBox Web, предназначенное для централизованного учета конфигураций, диагностических данных, файлов и сессий синхронизации ПЛК LogicBox в системах АСУТП. Система ориентирована на сопровождение инженерных данных и не выполняет удаленное управление промышленным оборудованием, что снижает риски эксплуатации и соответствует принятому ограничению проекта.

В первой главе была проанализирована предметная область и определены основные проблемы сопровождения ПЛК: разрозненное хранение конфигураций, отсутствие единого реестра контроллеров, сложность анализа диагностических данных и необходимость разграничения доступа. Во второй главе были спроектированы роли пользователей, функциональные требования, сценарии работы, модель данных, структура базы, архитектура приложения и механизм взаимодействия с локальным клиентом. В третьей главе была выполнена программная реализация системы на языке Go с использованием Fiber и PostgreSQL.

В результате разработки реализованы аутентификация и авторизация пользователей, ролевая модель доступа, разграничение прав по объектам эксплуатации, иерархия объектов вида организация --- участок --- стенд --- шкаф, привязка ПЛК только к шкафам, CRUD-интерфейсы для основных сущностей, загрузка и скачивание файлов, хранение и подтверждение конфигураций, просмотр диагностических записей, архивирование ПЛК, прием данных от локального клиента через API и журнал аудита действий пользователей.

Отдельное внимание было уделено целостности и безопасности данных. Критичные операции проверяются не только в интерфейсе, но и на уровне серверной логики и базы данных. Система контролирует допустимость иерархии объектов, блокирует некорректное удаление связанных данных, проверяет права пользователя перед выполнением операций, защищает обычные изменяющие запросы CSRF-токеном, а для API локального клиента использует Bearer-токен. Файлы хранятся в управляемом каталоге, а в базе данных фиксируются их метаданные и контрольные суммы.

Демонстрационная версия веб-приложения размещена по адресу \href{https://logicbox-web-halfdirty.waw0.amvera.tech/}{logicbox-web-halfdirty.waw0.amvera.tech}. Репозиторий с исходным кодом системы и LaTeX-материалами курсового проекта доступен по адресу \href{https://github.com/Half-dirty/logicbox-plc-webapp}{github.com/Half-dirty/logicbox-plc-webapp}.

Таким образом, цель курсового проекта достигнута. Разработанное веб-приложение обеспечивает централизованный учет ПЛК LogicBox, хранение конфигурационных и диагностических данных, работу с файлами, разграничение доступа, аудит действий и прием данных от локального клиента. Полученный результат соответствует тематике дисциплины «Разработка веб-приложений» и может использоваться как основа для дальнейшего развития системы сопровождения объектов АСУТП.
